% ----------------------------------------------------------
\chapter{Código de Normalização e Higienização de Dados}
\label{ap:normalizacao_coleta}
% ----------------------------------------------------------

Este apêndice apresenta o script desenvolvido em Python para carregar, extrair e normalizar os logs telemétricos das controladoras Wi-Fi Ruckus e UniFi. Devido às diferenças nativas nos formatos de reporte e nomes das chaves nos payloads JSON, o script realiza a consolidação das colunas comuns.

Abaixo é apresentado o trecho principal da rotina utilizada para extrair os logs brutos e mapear os dados para um formato CSV unificado de telemetria.

\begin{lstlisting}[language=Python, caption={Script de Normalização e Higienização da Rede UniFi}, label={lst:filtragem_unifi}]
import pandas as pd
import json

# Carregar o ficheiro CSV original contendo o log bruto
df = pd.read_csv('Logs-A-data-2026-06-17.csv')

# Funcao para transformar a string JSON da coluna 'Line' em dicionario
def parse_line(line):
    try:
        return json.loads(line)
    except:
        return {}

parsed = pd.json_normalize(df['Line'].apply(parse_line))

# Criar dicionario mapeando o MAC do AP para o seu Nome
ap_map = {}
if 'tipo' in parsed.columns and 'mac' in parsed.columns and 'name' in parsed.columns:
    ap_df = parsed[parsed['tipo'] == 'ap']
    for _, row in ap_df.iterrows():
        if pd.notna(row['mac']) and pd.notna(row['name']):
            ap_map[row['mac']] = row['name']

telemetria_rows = []

# Extracao de dados para formatar os registros
for _, row in parsed.iterrows():
    if row['tipo'] == 'cliente':
        ts = pd.to_datetime(row['ts']).strftime('%Y-%m-%d %H:%M')
        ap_mac = row.get('ap_mac', '')
        ap_name = ap_map.get(ap_mac, ap_mac)
        client_mac = row.get('mac', '')
        
        rssi = row.get('rssi', '')
        signal = row.get('signal', '')
        noise = row.get('noise', '')
        canal = row.get('channel', '')
        
        # Conversao de taxas brutas para Mbps
        tx_bps = float(row.get('tx_bytes_r', 0)) * 8
        rx_bps = float(row.get('rx_bytes_r', 0)) * 8
        tx_mbps = tx_bps / 1000000 
        rx_mbps = rx_bps / 1000000 
        
        # Modulacao fisica do Link Speed
        tx_rate_kbps = float(row.get('tx_rate', 0))
        rx_rate_kbps = float(row.get('rx_rate', 0))
        tx_link_mbps = tx_rate_kbps / 1000
        rx_link_mbps = rx_rate_kbps / 1000
        
        ccq = float(row.get('ccq', 1000))
        retries = (100 - (ccq / 10.0)) if pd.notna(ccq) else 0.0
        
        proto = row.get('radio_proto', '')
        padrao = {'ng': '802.11n', 'ac': '802.11ac', 'ax': '802.11ax', 'g': '802.11g'}.get(proto, proto)
        
        telemetria_rows.append({
            'timestamp': ts,
            'AP': ap_name,
            'cliente': client_mac,
            'RSSI': rssi,
            'signal': signal,
            'noise': noise,
            'canal': canal,
            'throughput_tx_mbps': tx_mbps,
            'throughput_rx_mbps': rx_mbps,
            'link_speed_tx_mbps': tx_link_mbps,
            'link_speed_rx_mbps': rx_link_mbps,
            'retransmissoes_pct': retries,
            'padrao': padrao
        })

df_output = pd.DataFrame(telemetria_rows)
df_output.to_csv('unifi_standardized.csv', index=False)
\end{lstlisting}
